% ============================================================
% Chapter 5: Development
%   Section: Fact-Checking Implementation
% Included from chapters/05-development/development.tex
% Sources: docs/decisions/incidents.md, retrieval.md, concurrency.md;
% docs/roadmap.md (Sep 21 audit, Sep 28, Sep 30, Oct 2).
% ============================================================

\section{Fact-Checking Implementation}
\label{sec:dev-fact-checking}

The design of Section~\ref{sec:design-fact-checking} is the result of
incidents as much as of plans. Each one below was reproduced, fixed and
pinned by a test, and each changed the design.

\paragraph{A re-analysed article was a duplicate of itself.} Every
extraction gives the article a new identifier, and the duplicate check
skipped a stored match only when the identifiers were equal. A second
analysis of the same URL --- a forced refresh, or any changed
threshold, since thresholds are part of the cache key --- found its own
earlier copy at similarity 1.0 and was rejected. Its quieter twin: the
corpus search had no exclusion at all, so the earlier copy came back as
evidence for its own claims. Both now exclude the article's URL inside
the vector query, and the store keeps one point per URL.

\paragraph{The title was never extracted.} trafilatura's second major
version leaves the title, author and date out unless asked for them.
Every article in the lake had no title, so everything built on the
headline --- restoring a claim's lost subject among it --- was dead. The
strategy's tests had mocked the library, which is why nobody noticed;
one test now runs it for real.

\paragraph{A claim was checked outside its article.} A Spanish sentence
about August prize contests, from an article about a street game, came
back \texttt{TRUE} at 84\% confirmed by the Mexico City marathon's prize
money. The fix restores the article's subject to the claim's queries
and gives the model the headline and the opening, with the rule that
evidence about a different event is unrelated even when it shares a
place, a date or a figure.

\paragraph{No call to the model without evidence.} Below the evidence
minimum the verdict was already forced to \texttt{UNVERIFIED}, yet the
model was still called: on 25~September, 23 of 24 claims spent the
slowest step of the pipeline on an answer that was thrown away. The
call is now skipped, and so is the progress event that told the
interface the model was being asked.

\paragraph{The model was retried behind the system's back.} The client
retried once, and the SDK underneath retried timeouts twice more by
default, so a slow local model could hold one claim for six times the
timeout and then report it unreachable. The SDK's own retries are off,
and a provider that was never reached raises an error that becomes the
claim's \texttt{llmUnreachable} flag, instead of returning nothing as a
malformed answer would.

\paragraph{The model's judgement of a source is enforced.} The verifier
already asked the model, per source, whether it supports, contradicts
or is unrelated to the claim, but a citation of a source the model had
just called unrelated still counted. Since 30~September it does not: a
definitive verdict resting only on such sources falls to
\texttt{UNVERIFIED} with a note saying why, at no extra model call.

\paragraph{Guardrails tested against a bluffing model.} A fake model ---
raw text behind the real client --- was made to bluff through the real
verifier and recalibration in 83 ways. Six got through and were fixed,
each with a test that failed before the fix: repeated citations
inflated the confidence (a one-source \texttt{FALSE} went from 0.70 to
1.00); a citation of \texttt{true} cited source~1; a null citation
list, or valid JSON that was not an object, failed the whole article; a
quote stitched from a headline and its body passed as verbatim; and a
per-run evidence minimum of zero sent a claim with no evidence to the
model.

\paragraph{Search failures look like verdicts.} The most consequential
finding of the project is not in the code: on 25~September, 68 of 69
search queries came back empty, because the engines behind SearXNG were
rate-limiting, challenging or timing out, and every affected claim came
back \texttt{UNVERIFIED}, indistinguishable from a real one. The
response was a closed engine list measured one by one (dead engines
removed, four academic indexes added at half weight, a three-second
timeout), half as many concurrent searches, each repeated query sent
once with non-empty answers cached for ten minutes, the
\texttt{searchUnavailable} state, and a keyless second route. On
1~October the cause of part of the rate-limiting was found in the
stack's own health check, which had sent a real search to every engine
every five seconds since September --- about 17,000 queries a day per
engine from an idle stack. It now asks the engine's own health
endpoint. In three runs that day every search answered.
